%% kapa: exam=1 round=32 year=2021 period=1 form=A subject=real_estate subject_ko=부동산학원론 seq=18 number=98
%% source: Q-Net 감정평가사 기출문제 (cst00309) artlSeq=5208397 file=extracted/1008/5208397/A형/2021년도 제32회 감정평가사 1차 1교시 A형.pdf
%% conversion: kordoc 4.24.51 -> tools/gen1.py (auto); review=proofread-sonnet; flags=none
\begin{question}{98}
부동산투자분석기법에 관한 설명으로 옳은 것을 모두 고른 것은? (단, 다른 조건은 동일함)
\begin{framedtext}
ㄱ. 현금유출의 현가합이 4천만원이고 현금유입의 현가합이 5천만원이라면, 수익성지수는 0.8이다. \par
ㄴ. 내부수익률은 투자로부터 발생하는 현재와 미래 현금흐름의 순현재가치를 1로 만드는 할인율을 말한다. \par
ㄷ. 재투자율로 내부수익률법에서는 요구수익률을 사용하지만, 순현재가치법에서는 시장이자율을 사용한다. \par
ㄹ. 내부수익률법, 순현재가치법, 수익성지수법은 할인현금흐름기법에 해당한다. \par
ㅁ. 내부수익률법에서는 내부수익률과 요구수익률을 비교하여 투자여부를 결정한다. \par
\end{framedtext}
\begin{choices}
  \item ㄱ, ㄹ
  \item ㄴ, ㄷ
  \item ㄹ, ㅁ
  \item ㄱ, ㄴ, ㅁ
  \item ㄷ, ㄹ, ㅁ
\end{choices}
\end{question}
\answer{3}
